\documentclass[11pt]{article}

\usepackage[margin=1in]{geometry}
\usepackage{amsmath,amssymb,amsthm}
\usepackage{enumitem}
\usepackage{xcolor}
\usepackage{mdframed}
\usepackage[colorlinks=true,linkcolor=black,urlcolor=blue!60!black,citecolor=black]{hyperref}
\usepackage{parskip}

\newcommand{\HH}{\mathbb{H}}
\newcommand{\QQ}{\mathbb{Q}}
\newcommand{\CC}{\mathbb{C}}
\newcommand{\ZZ}{\mathbb{Z}}
\newcommand{\tk}{\widetilde{k}_T}

\newmdenv[
  linecolor=gray!50,
  linewidth=0.5pt,
  backgroundcolor=gray!6,
  innertopmargin=8pt,
  innerbottommargin=8pt,
  innerleftmargin=10pt,
  innerrightmargin=10pt,
  roundcorner=3pt
]{promptbox}

\title{\vspace{-2.5em}Build brief --- the spine\\[0.3em]
\large DR--B periods \& \texorpdfstring{$L$}{L}-functions note}
\author{}
\date{}

\begin{document}
\maketitle
\vspace{-3.5em}

\section*{Contribution, stated narrowly}

The object is a finite, regulator-free, two-segment contour functional with a
Hurwitz-zeta $T$-kernel that delivers period polynomials and $L$-values from one
construction, recovers BFK on specialization, and extends --- within a homotopy
class --- to interior-pole meromorphic forms. The framework is Brown's
(arXiv:1710.07912): finite base point, $\tau_0$-independent cohomology class,
$T$-gauge. The contribution is the explicit closed kernel and the unified finite
realization, \emph{not} new objects --- the $L$-values are BFK's, the periods are
Brown's quasi-periods made closed-form.

Language: a manifestly convergent finite formula that \emph{equals} the regularized
value is a certificate of canonicality --- scheme-, continuation-, and
cusp-choice-independent --- not a correction to results that were already rigorous.
Lean on \emph{canonical / certified}, never on \emph{correct}.

\section*{The spine}

Lead with the construction question, so the kernel is an \emph{answer}:

\begin{center}\itshape
Periods are a cocycle. How do you build one --- parabolic --- for a modular form
from a generic finite base point?
\end{center}

\paragraph{Precision 1: the gate is the constant term, not the pole structure.}
The parabolic gauge (killing $C_T$) exists iff $c_f(0)=0$; the principal part never
obstructs it, interior poles included. ``Generic pole structure'' is the right
phrase and a selling point --- arbitrary poles are fine. But ``generic form'' at the
cocycle step is too generous: $c_f(0)\neq0$ (an Eisenstein series) cannot be made
parabolic and yields the rational/Eisenstein piece. Clean statement: \emph{any form
with vanishing constant term, regardless of pole structure, admits the parabolic
gauge.} The gate is one Fourier coefficient.

\paragraph{Precision 2: the two legs have different gates --- don't bundle them.}
The period-polynomial leg needs the parabolic gauge ($c_f(0)=0$, i.e.\ $f\in S_k^!$).
The $L$-function leg does not: the two-segment functional is defined on all of
$M_k^!$, constant term allowed, with the expected poles of $\Lambda$ coming from
$c_f(0)$ as classically. The arc \emph{widens} as it goes; let that be visible.

\paragraph{Presentation: define and verify; derive in a footnote.}
Do not walk the reader through the variation-in-$T$ operator, the linear algebra,
and the Bernoulli combinatorics. Present $\tk$ as a definition, apply $\delta_T$, and
show the parabolic obstruction vanishes (for $f$ with $c_f(0)=0$) at all $s$ --- a
two-line check. Push the provenance to a footnote: (1) construct the
variation-in-$T$ operator $\delta_T$; (2) set up the coboundary linear algebra;
(3) read off the combinatorics; (4) recognize them as Bernoulli, i.e.\ special
values of $\zeta(s,x)$ at integer $s$; (5) note that this $\zeta(s,x)$ is already
entire in $x$ and meromorphic in $s$, so the inversion is a \emph{function} identity
with the integers as one slice. The footnote says the integer case \emph{revealed
which analytic object was operating} --- not that a discrete identity was
analytically continued. The second framing invites a continuation-validity question
the construction does not carry.

\paragraph{The four beats ($\sim$5 pp with disciplined front/back matter).}
\begin{enumerate}[leftmargin=1.6em,topsep=2pt,itemsep=4pt]
\item \textbf{Periods are a cocycle; the construction question.} Eichler--Shimura in
a paragraph; the task is a parabolic cocycle from a generic finite base point. This
is Brown (arXiv:1710.07912) --- recall it, with the gate $c_f(0)=0$ and its earlier
form as constant-term subtraction $\widetilde f|(1-T)=0$ (arXiv:1202.5802).
\item \textbf{Define the Hurwitz $T$-kernel; verify it inverts $\delta_T$.}
$\tk(\tau,s)=\zeta(1-s,\tau+1)-e^{i\pi(s-1)}\zeta(s-(k-1),\tau+1)$, implicit in
Brown's $P$; the contribution is the explicit all-$s$ form. Verify, don't derive:
$\delta_T\,\zeta(1-s,\tau+1)=-(\tau+1)^{s-1}$ inverts $\delta_T$ uniformly in $s$, so
applying $\delta_T$ to the kernel leaves exactly the constant-term piece --- which is
why the obstruction vanishes \emph{iff} $c_f(0)=0$. State that gate here even in the
short version; the verification is precisely where it appears. On $S_k^!$ this beat
delivers the period polynomial $r_f$/$\omega^\pm$ as the two-segment functional.
\item \textbf{Swap the period monomial for $\tau^{s-1}$ $\to$ the $L$-function.}
Same kernel; the functional is entire in $s$, and $L^*(f,s)$ is entire on $S_k^!$,
with the expected poles otherwise --- on $M_k^!$ it is meromorphic, carrying exactly
the poles forced by $c_f(0)$, as classical $\Lambda$ does. (Say ``entire on $S_k^!$,
expected poles otherwise''; do not call the $M_k^!$ object entire.) The $\tau_0=i$
specialization recovers BFK with the pair $\Gamma(s,2\pi n)$ and
$i^k\Gamma(k-s,2\pi n)$ --- the $s\leftrightarrow k-s$ reflection making
$\Lambda^*(s)=i^k\Lambda^*(k-s)$ manifest; $b=1$ is the split \emph{height}
($\tau=i$, the $S$-fixed point), not the $s$-center. All-$s$ is a \emph{result}: the
term-by-term unfolding gives a manifestly meromorphic incomplete-gamma series
matching the classical / Apostol / BFK expressions coefficient by coefficient, so
there is no continuation step --- only the finite principal-part interchange (the
Hurwitz-unfolding proposition).
\item \textbf{Encapsulation.} One object covers both literature regimes: weak /
regular cusps with arbitrary principal part (fully canonical, simply connected),
\emph{and} interior-pole meromorphic forms --- canonical given a contour, the
cross-wall dependence being residues. State the winding prescription explicitly (the
symmetric $\tfrac12\sum_{\pm\epsilon}$ of arXiv:1806.09874): the identity is then
definitional. The jump at an interior pole is exactly $f$'s residue times the finite
kernel value there --- $\zeta(\sigma,\tau+1)$ is holomorphic throughout $\HH$, so the
kernel contributes no residue of its own.
\end{enumerate}

\section*{References to genuflect toward}

BFK (2014); arXiv:1809.02938 (the incomplete-gamma $L$ the unfolding reproduces);
arXiv:1603.03056; arXiv:2202.09542; Diamantis et al.\ (RMS 2020); Brown--Fonseca
(arXiv:2508.04844). Frame the contribution as the finite, regulator-free
Hurwitz-kernel \emph{representation} unifying periods and $L$-values.

\section*{Paste-ready prompt for a Claude Code instance}

\begin{promptbox}
\small
Write a $\sim$5-page note from this spine. Define-and-verify, not derive-and-discover.
\begin{enumerate}[leftmargin=1.4em,topsep=1pt,itemsep=1pt,nosep]
\item Open with the construction question: how to build a parabolic cocycle for a
modular form from a generic finite base point. Recall Brown (arXiv:1710.07912) in one
paragraph; state the gate $c_f(0)=0$.
\item Present the Hurwitz $T$-kernel
$\tk(\tau,s)=\zeta(1-s,\tau+1)-e^{i\pi(s-1)}\zeta(s-(k-1),\tau+1)$ as a definition.
Verify $\delta_T\,\zeta(1-s,\tau+1)=-(\tau+1)^{s-1}$ and that the parabolic obstruction
vanishes iff $c_f(0)=0$. Footnote the five-step provenance; frame step 5 as revealing
the entire-in-$x$, meromorphic-in-$s$ object, not as continuing a discrete identity.
\item Period polynomials on $S_k^!$ from the two-segment functional.
\item $L$-function by swapping the period monomial for $\tau^{s-1}$: entire on $S_k^!$,
expected poles on $M_k^!$. $\tau_0=i$ recovers BFK with $\Gamma(s,2\pi n)$,
$i^k\Gamma(k-s,2\pi n)$; $b=1$ is the split height, not the $s$-center. State all-$s$
as a result (coefficient-by-coefficient match to the incomplete-gamma series), not a
sketch.
\item Encapsulation: weak/regular cusps (fully canonical) and interior poles
(canonical given a contour; cross-wall jumps are $f$'s residues times the finite
kernel value). State the symmetric winding prescription explicitly.
\item Relate to prior work: BFK (2014), arXiv:1809.02938, 1603.03056, 2202.09542,
Diamantis et al.\ (RMS 2020), Brown--Fonseca (arXiv:2508.04844). Contribution is the
finite regulator-free Hurwitz-kernel representation, not new objects.
\end{enumerate}
Tone: canonical/certified, never ``correct''; parsimonious novelty language; no claim
broader than the simply-connected case without the winding clause.
\end{promptbox}

\end{document}
